
est aussi repr\'esentable par une immersion ferm\'ee.}
\item \emph{\(\underline{\Pic}^\tau_{X/S}\) est de type fini sur \(S\), c'est-\`a-dire, la famille des classes des faisceaux inversibles \(\tau\)-\'equivalents \`a \(0\) sur les fibres de \(X/S\) est limit\'ee.}
\end{enumerate}

En effet, dans (i) (resp. (ii)), on doit traiter le cas d'un morphisme \(T\to\underline{\Pic}_{X/S}\), o\`u \(T\) est un \(S\)-sch\'ema quelconque, et montrer que l'image inverse de \(\underline{\Pic}^\tau_{X/S}\) est repr\'esentable par un sous-sch\'ema ouvert (resp. et ferm\'e) de \(T\). Il suffit de prendre \(T=\Spec(B)\) sur un ouvert affine \(\Spec(A)\) de \(S\). Alors \(T\) est la limite du syst\`eme projectif filtrant des \(S\)-sch\'emas affines \(T_\alpha=\Spec(B_\alpha)\), o\`u \(B_\alpha\) parcourt les sous-\(A\)-alg\`ebres de type fini de \(B\). En vertu de 3.1, le morphisme en question provient d'un morphisme \(T_\alpha\to\underline{\Pic}_{X/S}\). On peut donc supposer \(T\) de type fini sur \(S\).

Le morphisme \(T\to\underline{\Pic}_{X/S}\) correspond \`a un morphisme f.p.p.f. \(R\to T\) et un faisceau inversible \(L\) sur \(X_R\). Par descente f.p.p.f., on peut remplacer \(T\) par \(R\). Apr\`es le changement de base \(R\to S\), compte tenu de la compatibilit\'e de la formation de \(\underline{\Pic}^\tau_{X/S}\) avec tout changement de base (4.1) on est donc ramen\'e dans (i) et (ii) \`a \'etudier le cas d'une section \(S\to\underline{\Pic}_{X/S}\), d\'efinie par un faisceau inversible \(L\) sur \(X\).

On doit prouver que l'ensemble \(S_0\) des \(s\in S\) tels que \(L_s\) soit \(\tau\)-\'equivalent \`a \(0\) est ouvert, et aussi ferm\'e sous les conditions de (ii)\,; comme \(S\) est noeth\'erien, ceci veut dire trois choses~: pour un \(s\in S_0\),
\begin{enumerate}[label=(\alph*)]
\item toute g\'en\'erisation de \(s\) reste dans \(S_0\).
\item toute sp\'ecialisation de \(s\), voisine de \(s\), reste dans \(S_0\).
\item sous les conditions de (ii), toute sp\'ecialisation de \(s\) reste dans \(S_0\).
\end{enumerate}

